\documentclass[11pt]{article}
\usepackage[margin=1in]{geometry}
\usepackage{amsmath,amssymb,amsthm,mathtools,booktabs,array}
\usepackage[T1]{fontenc}
\usepackage{lmodern,microtype}
\usepackage{xcolor}
\usepackage[colorlinks=true,linkcolor=blue!45!black,citecolor=blue!45!black,urlcolor=blue!45!black]{hyperref}
\usepackage{enumitem}
\setlist{itemsep=3pt,topsep=5pt}
\setlength{\emergencystretch}{2em}
\newcommand{\C}{\mathbb C}
\newcommand{\R}{\mathbb R}
\newcommand{\N}{\mathbb N}
\DeclareMathOperator{\Tr}{Tr}
\DeclareMathOperator{\Res}{Res}
\newtheorem{theorem}{Theorem}[section]
\newtheorem{proposition}[theorem]{Proposition}
\newtheorem{lemma}[theorem]{Lemma}
\newtheorem{corollary}[theorem]{Corollary}
\theoremstyle{remark}
\newtheorem{remark}[theorem]{Remark}
\hypersetup{pdftitle={Complete-source transport, conserved contours and pair-spectrum bounds for the theta heat flow},pdfauthor={Drafted for Edward Baker}}
\title{Complete-source transport, conserved contours and pair-spectrum bounds for the theta heat flow}
\author{Drafted for Edward Baker}
\date{Working research draft\\23 October 2026}
\begin{document}
\maketitle
\begin{abstract}
We develop finite-population estimates for the canonical theta heat flow,
retaining the complete field and its exterior contribution. The bounded
pair-spectrum trace admits an exact endpoint collision value, a sharp
finite multiplicity consumer, and a forest-polynomial representation in
prepared coefficients. A contour equation for those coefficients is driven
by commuting degree-lowering operators; at fixed degree its backward
propagator is an explicit finite polynomial in integrated source moments.
These statements permit quantitative error accounting without assuming
internal root separation. We also construct a holomorphic correction to
the effective Riemann--Siegel reference and give its field and heat-defect
budgets in a microscopic real-time regime. A complete-source interval
calculation on one fixed low-height box certifies its entire two-root
population and discriminant throughout a complex time square. Direct
evaluation is stronger for the terminal filter in this comparison. A bounded
complete-cap calculation also encloses every face, but both lower bounds
are negative and the direct interval is slightly narrower.
A provisional extension develops conservation of the heat one-form on
moving spatial contours, primitive and adjoint Noether actions, and an
infinite hierarchy of Fourier and polynomial invariants. Explicit
collision controls distinguish these exact laws from the still missing
canonical deformation and one-sided flux estimates.
The height-uniform microscopic upper estimate required for a global
simplicity conclusion remains unproved; this manuscript proves neither
RH nor a new global simplicity proportion.
\end{abstract}

\section{Purpose, scope and normalization}
The aim is to turn complete source information into quantitative control
of a finite collection of zeros of the theta heat flow, and to determine
which resulting bounds could support a global simplicity theorem. We
separate three levels of result: finite identities with paid perturbation
bounds; a bounded computation for the unchanged theta field; and the
independent, still missing source estimate needed at growing height.
This distinction is essential when a generic heat identity also holds
for collision controls.

Throughout,
\begin{align}
H(\tau,z)&=\int_0^\infty e^{\tau u^2}\Phi(u)\cos(zu)\,du,
\label{eq:canonical}\\
\Phi(u)&=\sum_{n\ge1}
 (2\pi^2n^4e^{9u}-3\pi n^2e^{5u})e^{-\pi n^2e^{4u}}.
\nonumber
\end{align}
The source is positive for \(u\ge0\), has superexponential decay, and defines
an entire function of both variables. With the standard completed zeta
function,
\begin{equation}
 H_\tau=-H_{zz},\qquad 8H(0,z)=\xi(1/2+iz/2).
 \label{eq:normalization}
\end{equation}
These conventions agree with \cite{polymath,rodgers-tao}; the real
\(H\)-coordinate is twice the usual zeta ordinate.

A \emph{complete finite population} means every spatial zero inside a
declared enclosing wall, counted with multiplicity. A nonzero wall keeps
this count constant, including through internal collisions. Its monic
prepared polynomial \(P\) gives \(H=UP\) in the interior, with \(U\)
nonvanishing. The factor \(U\), equivalently complete field data on the
wall, retains the exterior contribution. None of the results below
permits replacing this population by a selected subset of convenient
real roots.

The main results are organized by their logical roles.
The spectral and forest results give consumers of source estimates.
The coefficient equation provides a finite transport mechanism.
The corrected reference supplies explicit source and defect budgets,
but still requires a whole-wall lower bound. The finite theta certificate
shows that the interfaces can actually be supplied on one prescribed
box. The physical-cap comparison retains the full boundary and its
spatial complement. The provisional spatial-conservation section supplies
a separate analytic direction whose useful canonical estimates remain open.

\paragraph{Assumptions and research status.}
Real-spectrum positivity and the global adjacent-gap comparison are used
only where explicitly stated. A local Hermite birth, a complex contour
identity, or a finite complete-source certificate does not itself require
RH. Conversely, excluding multiple roots does not establish reality of
all roots and therefore does not prove RH.
No priority claim is made for the finite algebraic identities or the
combination of classical ingredients. The repository research notes give
the derivation history; this draft collects selected statements and
proofs into one reviewable account.

% Prepared for Edward Baker, 23 October 2026.
% Substantial LLM assistance: gpt-6-astra (Codex), reasoning effort ultra.
% Working manuscript contribution; auxiliary details are in Notes 10 and 11.
% Expository clarification, 23 October 2026: gpt-6.1-sol (Codex), effort xhigh.

\section{Collision spectra, positive-time transfer, and the missing global estimate}
\label{sec:spectral-transfer}

For distinct real points $x_1,\ldots,x_N$, let $e_i$ denote the $i$th
standard basis column vector in $\R^N$, and write
\[
 L(x)=\sum_{i<j}\frac{(e_i-e_j)(e_i-e_j)^T}{(x_i-x_j)^2},
 \qquad q(t)=\Tr\!\left(4tL(x(t))[I+4tL(x(t))]^{-1}\right).
\]

\paragraph{Geometric meaning and choice of the filter.}
Let $V(x)=(x_i^{j-1})_{i,j=1}^N$ be the Vandermonde matrix and put
$\Delta=(\det V)^2=\prod_{i<j}(x_i-x_j)^2$. Direct differentiation gives
\begin{equation}
 L=-\operatorname{Hess}_x\log|\det V|
   =-\tfrac12\operatorname{Hess}_x\log\Delta,
 \qquad
 v^TLv=\sum_{i<j}\frac{(v_i-v_j)^2}{(x_i-x_j)^2}.
 \label{eq:vandermonde-hessian}
\end{equation}
Thus $L$ is the curvature matrix of the logarithmic repulsion energy,
and also the weighted graph Laplacian with inverse-square gap weights.
The Hessian here is the full matrix of second derivatives with respect to
the root positions; the scalar spatial Laplacian would give its trace.
The constant vector is a zero mode because a common translation changes
none of the gaps.

For $t>0$ on a real simple slice, let $\lambda_k(t)\ge0$ be the eigenvalues
of $L(x(t))$. Then
\begin{equation}
 q(t)=\sum_{k=1}^N f(4t\lambda_k(t)),
 \qquad f(u)=\frac{u}{1+u}.
 \label{eq:filter-effective-dimension}
\end{equation}
This is a smooth count, or effective dimension, of the modes at the heat
scale: a mode with $4t\lambda_k\ll1$ contributes little, a mode with
$4t\lambda_k\gg1$ contributes nearly one, and a mode at $4t\lambda_k=1$
contributes one half. Each mode contributes at most one, so a single
exceptionally small gap cannot give an arbitrarily large contribution
through one eigenvalue.

The factor $4t$ is calibrated by the elementary heat birth
$H(t,z)=z^2-2t$. Its roots are $\pm\sqrt{2t}$, their squared gap is $8t$,
and their two-point Laplacian has nonzero eigenvalue $2/(8t)=1/(4t)$.
This double therefore contributes exactly $1/2$ to $q$.
The proposition below gives the corresponding spectrum for every local
multiplicity. For a fixed finite population of distinct real endpoint
roots, by contrast, $L(x(t))$ stays bounded and $q(t)\to0$ as $t\downarrow0$.
At positive time, very close descendants of distinct endpoint roots can
also contribute; $q$ is therefore used to bound multiplicity.

The particular function $f$ is a convenient choice of smooth filter.
Besides boundedness, it has the matrix order property
$f(A)=I-(I+A)^{-1}$ for $A\succeq0$, used in the comparison proof below,
and a determinant representation:
\begin{equation}
 q(t)=\left.
 a\,\partial_a\log\det\!\left(I+aL(x(t))\right)
 \right|_{a=4t}.
 \label{eq:filter-log-determinant}
\end{equation}
Here the root configuration and time are held fixed in $\partial_a$;
this is not the total derivative along the heat flow.
Thus the filter is also the logarithmic scale derivative of the determinant
of the regularized curvature matrix. This identity leads to the
forest-polynomial representation in Section~\ref{sec:forest-filter}.

A selected population is \emph{complete at its endpoint sites} if it includes
all multiplicity labels at every site it selects. Its finite matrix contains
only those labels, but their motion is always that of the complete field.
For endpoint multiplicities $m_a$, distinguish
\[
 M_0=\sum_{m_a\ge2}m_a,\qquad D_0=\sum_a(m_a-1),\qquad N=\sum_a m_a.
\]
Thus $M_0$ counts labels at nonsimple sites; $D_0$ counts excess multiplicity.

\begin{proposition}[Endpoint spectrum and scalar multiplicity encoding]
\label{prop:endpoint-spectrum}
Let a finite complete endpoint population of a real-type analytic solution of
$H_t=-H_{zz}$ have real sites, with real simple descendants for small $t>0$.
Then
\[
 \operatorname{spec}(4tL(x(t)))\longrightarrow
 \bigsqcup_a\{0,1,\ldots,m_a-1\},
 \qquad
 q_0:=\lim_{t\downarrow0}q(t)=\sum_a(m_a-\mathsf H_{m_a}),
 \quad \mathsf H_m=\sum_{j=1}^m\frac1j.
\]
In particular $D_0/2\le q_0\le D_0$ and $M_0\le4q_0$.
\end{proposition}
\begin{proof}
The local heat expansion gives
$x_{a,j}(t)=a+\sqrt{2t}\,h_{m_a,j}+O(t)$, where the $h_{m,j}$ are the
zeros of the probabilists' monic Hermite polynomial $\mathrm{He}_m$.
The Hermite equation implies
$\sum_{j\ne i}(h_i-h_j)^{-1}=h_i/2$.
Acting with $L(h)$ on evaluations of the monomial $x^k$ gives
$(k/2)h_i^k$ plus lower-degree terms. The evaluations of degrees
$0,\ldots,m-1$ form a Vandermonde basis, so
$\operatorname{spec}L(h)=\{0,1/2,\ldots,(m-1)/2\}$.
Inter-site entries, including their diagonal contributions, vanish after
multiplication by $t$. The block limits therefore give the stated spectrum.
Applying $u\mapsto u/(1+u)$ yields
$\sum_{k=1}^{m-1}k/(k+1)=m-\mathsf H_m$.
Each positive summand lies between $1/2$ and $1$; also
$m-\mathsf H_m\ge m/4$ for $m\ge2$, with equality at $m=2$.
\end{proof}

The Hermite spectrum is inherited from the local analysis in Note 06; its
scalar trace is an exact encoding, not a theta upper estimate. A double
contributes $1/2$, a triple $7/6$, and a quadruple $23/12$.
These are fixed-population limits, with no uniform exchange of height and
time limits asserted.

\begin{theorem}[Sharp positive-time transfer]
\label{thm:sharp-transfer}
Suppose, in addition, that all adjacent descendants within each selected
endpoint cluster remain globally adjacent in a complete real simple system
on $(0,T]$, and that their gaps satisfy
\[
 (g^2)'\le8,\qquad g(0)=0.
\]
Then $M_0\le4q(t)$ for every $0<t\le T$. The constant four is sharp.
\end{theorem}
\begin{proof}
Every such gap obeys $g(t)^2\le8t$. Partition each even cluster into adjacent
pairs, and each odd cluster of size at least three into pairs and one
consecutive triple. Retain only the pair edges and the two consecutive edges
in each triple. Then $4tL\succeq A_*$, where the pair and triple blocks of
$A_*$ are one half of the respective path Laplacians.
The map $f(A)=I-(I+A)^{-1}$ is operator monotone on positive semidefinite
matrices. A pair block has spectrum $\{0,1\}$ and contributes $1/2$ to
$\Tr f(A_*)$. A triple has spectrum $\{0,1/2,3/2\}$ and contributes
$1/3+3/5=14/15\ge3/4$. Thus each block pays at least one quarter of its
labels. Singleton sites need no payment. For sharpness,
$H(t,z)=z^2-2t$ gives $M_0=2$ and $q(t)=1/2$.
\end{proof}

For canonical theta, the hypothesis on gaps is supplied by the full
canonical real-root motion. For globally adjacent roots $a<b$, $g=b-a$,
\[
 (g^2)'=8-4g^2\sum_{r\ne a,b}\frac1{(a-r)(b-r)}\le8.
\]
The exterior sum is part of the hypothesis and is absolutely convergent;
its terms are nonnegative by global adjacency. No external factor is set to
one. In the conditional theta application, RH supplies the required complete
real system at positive time. The finite matrix need not contain the exterior
roots, but a truncated motion law is insufficient. The improvement from the
earlier constant six concerns proportional transfer; the sufficient
all-simple threshold remains $q(t)<1/2$.

\begin{theorem}[Bounded spectral mass and short adjacent gaps]
\label{thm:gap-filter}
For arbitrary $y_1<\cdots<y_N$, set $A=L(y)$ and
\[
 B_N(\delta)=\frac1N\#\{i<N:y_{i+1}-y_i<\delta\},\qquad
 Q_N(s)=\frac1N\Tr\!\left(sA(I+sA)^{-1}\right).
\]
For $s,\delta>0$,
\begin{equation}
 B_N(\delta)\frac{s}{\delta^2+2s}
 \le Q_N(s)\le2B_N(\delta)+\frac{\pi^2s}{3\delta^2}.
 \label{eq:gap-filter-two-sided}
\end{equation}
Consequently, along any sequence of configurations and positive scales,
\[
 Q_N(s_N)\longrightarrow0
 \quad\Longleftrightarrow\quad
 B_N(a\sqrt{s_N})\longrightarrow0\quad\hbox{for every fixed }a>0.
\]
\end{theorem}
\begin{proof}
Mark every vertex touching a short adjacent gap. If $r$ vertices are marked,
then $r\le2NB_N(\delta)$. The unmarked vertices are $\delta$-separated;
every marked point is also at distance at least $\delta$ from each
unmarked point. Otherwise the unmarked point's nearest original neighbor
on that side would be less than $\delta$ away.

Let $D$ be the principal submatrix of $A$ on the unmarked set. Its diagonal
\emph{retains} interactions with marked vertices. For an unmarked vertex,
its $k$th unmarked neighbor on either side is at distance at least $k\delta$.
For a marked vertex, its $k$th unmarked neighbor on either side satisfies the
same bound. Summing $2\sum_{k\ge1}(k\delta)^{-2}$ first over the unmarked
vertices and then over marked-to-unmarked interactions gives
\[
 \Tr D\le\frac{\pi^2N}{3\delta^2}.
\]
Interlacing pays at most $r$ for the discarded eigenvalues; the remaining
filter values are bounded by $s$ times the eigenvalues of $D$. This proves
the upper bound, including the empty-principal-matrix case.

The short adjacent edges form disjoint paths. From $K$ such edges select a
matching of at least $K/2$ edges. Each retained edge has a Laplacian
eigenvalue at least $2/\delta^2$. Dropping all other edges and using operator
monotonicity proves the lower bound. Substituting $\delta=a\sqrt{s}$ gives
\[
 B_N(a\sqrt{s})\le(a^2+2)Q_N(s),\qquad
 Q_N(s)\le2B_N(a\sqrt{s})+\frac{\pi^2}{3a^2}.
\]
The equivalence follows by taking a limsup at fixed $a$, then letting
$a\to\infty$ for the reverse direction.
\end{proof}

For example, independently acquired bounds
$B_N(\delta)\le C\delta^\beta+\varepsilon_N$, $\beta>0$, imply
\[
 Q_N(s)\le\left(2C+\frac{\pi^2}{3}\right)s^{\beta/(\beta+2)}
             +2\varepsilon_N
\]
whenever $s^{1/(\beta+2)}$ lies in their certified range. This is a sufficient
source input, not a bound proved here for theta.

To state the global interface precisely, let an endpoint target have $N_h$
labels, of which $n_h=N_h-r_h$ are covered by complete clusters. At a
source-prescribed time $t_h$, set $y_i=\nu_hx_i(t_h)$ and
$s_h=4t_h\nu_h^2$. Under Theorem~\ref{thm:sharp-transfer},
\begin{equation}
 \frac{M_h}{N_h}\le\frac{r_h}{N_h}
       +4\frac{n_h}{N_h}Q_{n_h}(s_h).
 \label{eq:global-spectral-interface}
\end{equation}
Thus density-one simplicity would follow conditionally from $r_h=o(N_h)$
and a vanishing fraction of gaps smaller than $2a\sqrt{t_h}$ for every fixed
$a>0$. That microscopic gap estimate remains unpaid. Choosing times after
resolving the unknown endpoint gaps would merely reformulate the target.
The published late-clock onset $x\ge e^{C/t}$ in \cite{polymath} does not
supply this estimate. With its natural unfolding
$\nu_X=g_x(X,t)\sim\log X/(4\pi)$, the uniformly bounded normalized gaps
instead imply $Q_N(s)\ge1-1/N-(1+4s/C_g^2)^{-1/2}$ by path comparison.
For this unfolding, $s=4t\nu_X^2\to\infty$ in the high-height onset regime,
so the unmodified filter saturates. This does not exclude
matched backward transport on a different clock.

\section{A coefficient-level filter with explicit conditioning}
\label{sec:forest-filter}

The preceding filter is defined through roots, but its acquisition need not
isolate them individually. Let $P(z)=\prod_{i=1}^N(z-x_i)$ be monic and
$\Delta=\prod_{i<j}(x_i-x_j)^2$. The following identity is algebraic; real
rootedness is needed only for the positivity conclusions.

\begin{proposition}[Forest polynomial and coefficient enclosure]
\label{prop:forest-filter}
For $a>0$ and distinct roots define $D=\Delta\det(I+aL)$ and $G=a\partial_aD$.
Then
\begin{equation}
 D=\sum_{F\,\operatorname{forest}}
     a^{|E(F)|}\prod_{C\in\operatorname{comp}(F)}|C|
     \prod_{\{i,j\}\notin E(F)}(x_i-x_j)^2,
 \qquad q_a=\frac GD.
 \label{eq:forest-polynomial}
\end{equation}
The sum runs over all spanning labeled forests, including isolated vertices.
The quotient identity is asserted where $D\ne0$. Both $D$ and $G$ are
symmetric polynomials in the roots and hence polynomials in the coefficients
of $P$. Fix $q_*=G_*/D_*$ and form the coupled polynomial $K=G-q_*D$.
If coefficient data give $|D-D_*|\le E_D<|D_*|$ and $|K|\le E_K$, then
\begin{equation}
 |q_a-q_*|\le\frac{E_K}{|D_*|-E_D}.
 \label{eq:forest-residual-bound}
\end{equation}
The separate allowance $E_K=\epsilon_G+|q_*|\epsilon_D$, obtained from
$|G-G_*|\le\epsilon_G$, is a valid fallback.
\end{proposition}
\begin{proof}
Apply Cauchy--Binet to the weighted incidence matrix in $L$.
A cyclic edge selection has dependent columns and contributes zero.
For a forest, its incidence Gram determinant is $\prod_C|C|$: its nonzero
maximal minors omit one vertex per component and have square one.
Thus its determinant contribution is the edge-weight product times
$a^{|E|}\prod_C|C|$. Each edge occurs at most once, so multiplying
by the single product $\Delta$ cancels every denominator. This proves the
polynomial identity. Logarithmic differentiation of the determinant gives
\[
 a\partial_a\log D=a\Tr\bigl(L(I+aL)^{-1}\bigr)=q_a.
\]
The discriminant is independent of $a$. Finally $q_a-q_*=K/D$, and
$|D|\ge|D_*|-E_D$, proving \eqref{eq:forest-residual-bound}.
\end{proof}

The coefficient residuals in the proposition are finite, computable quantities.
For example, on a coefficient polydisc $|c_j-c_{*,j}|\le r_j$, take
$D_*=D(c_*)$, $G_*=G(c_*)$ and use the finite Taylor majorants
$E_D=\sum_{\alpha\ne0}|\partial^\alpha D(c_*)|r^\alpha/\alpha!$
and
$E_K=\sum_{\alpha\ne0}|\partial^\alpha K(c_*)|r^\alpha/\alpha!$,
since $K(c_*)=0$. Form $K$ before taking absolute values, so correlated
coefficient changes can cancel. The estimating gain is that the paid
margin is for $D$, not necessarily for $\Delta$ or individual roots. It is
not a statement that all coefficient balls are well conditioned.

\begin{lemma}[Real repeated-root boundary]
For fixed $a>0$ and a real root multiset, the polynomial extension has $D>0$
if and only if every root multiplicity is at most two.
\end{lemma}
\begin{proof}
All forest terms are nonnegative. A nonzero term must include every edge
joining equal roots, since any omitted such edge contributes zero. A cluster
of three or more forces a cycle and makes every term vanish. If all clusters
have size at most two, the disjoint collapsed pair edges form a forest whose
complementary gaps are nonzero, so its term is positive.
\end{proof}

The distinction is visible already in degrees two and three. For
$P=z^2+bz+c$, put $\delta=b^2-4c$. Then
\[
 D=\delta+2a,\qquad G=2a,\qquad q_a=\frac{2a}{\delta+2a}.
\]
In particular, $P=z^2+p$, $-a/8\le p\le0$, gives $4/5\le q_a\le1$,
including the repeated reference polynomial. There is no discriminant floor
or internal root isolation in this enclosure. It is a stable filter
calculation, not a simplicity certificate: its values are above the relevant
threshold. For complex $|p|\le a/8$, one similarly has
$|D|\ge3a/2$ and $|q_a-1|\le1/3$.

For $P=z^3+bz^2+cz+d$, let $B=b^2-3c$. Directly from the three-vertex forest
sum,
\[
 D=\Delta+2aB^2+6a^2B,\qquad G=2aB^2+12a^2B.
\]
On the real-rooted cone, $0\le\Delta\le4B^3/27$. Thus, for $B>0$ and
$\kappa=B/a$,
\[
 \frac{2\kappa+12}{4\kappa^2/27+2\kappa+6}
 \le q_a\le\frac{2\kappa+12}{2\kappa+6}.
\]
The restricted real limit at a triple is two. There is no corresponding
universal extension in complex coefficient space: along $P=z^3-r^3$,
$B=0$, $\Delta=-27r^6$, and $q_a=0$, whereas the real Hermite approach gives
$q_a\to2$ at fixed $a$.

The heat scale must also be retained. For the cubic heat polynomial
$P_t=z^3-6tz$,
\[
 D(t,a)=108t(a+2t)(a+4t),\qquad
 D(t,4t)=5184t^3,\qquad G(t,4t)=6048t^3,
\]
so $q(t)=7/6$. More generally, the coefficient snapshots
\[
 P_t=z^3-6tz+\eta t^{3/2},\qquad
 q(t)=\frac{6048}{5184-27\eta^2}.
\]
For $\eta\ne0$ these snapshots do not form a heat solution.
For real $\eta$ with $|\eta|\le\sqrt{32}$ the polynomials are real rooted
and $7/6\le q\le7/5$. A constant coefficient perturbation of order $t^{3/2}$
therefore has an order-one effect on the limiting filter. To recover the
Hermite value, sufficient depressed-coefficient accuracies are
$p+6t=o(t)$ and $d_{\rm depressed}=o(t^{3/2})$; the cost of centering must
be included. Likewise the quadratic heat polynomial retains the scale
$\epsilon/t$ for its constant coefficient error. Exact coefficient transport
removes rootwise inverse-gap factors, but cannot remove these source precision
requirements. Notes 10 and 11 give further controls and the associated
complete finite replays.

\section{Complete-source coefficient transport on a fixed wall}
\label{sec:coefficient-transport}

The root velocity equations introduce inverse internal gaps. Passing to the
complete prepared polynomial removes those denominators from the transport
equation. The source information is still required, but it can be acquired on
one exterior wall. We give the finite theorem with its source cost explicit.

Fix a circle $\Gamma=\{|z-z_0|=R\}$, positively oriented, and a real time
interval $[t,T]$. Suppose $H$ is holomorphic on a neighborhood of the closed
spatial disk and the time interval, satisfies
\begin{equation}\label{eq:heat-fixed-wall}
 \partial_tH=-\partial_z^2H,
 \qquad H(u,w)\ne0\quad(u\in[t,T],\ w\in\Gamma),
\end{equation}
and has $N$ zeros inside $\Gamma$, counted with multiplicity. Write
$H=UP$, where $P$ is the complete monic polynomial of those zeros and $U$ is
holomorphic and nonzero in the disk. The center and radius remain fixed.
For a polynomial $f$, set
\[
 Q_f(w,z)=\frac{f(w)-f(z)}{w-z},\qquad
 Z_H=\frac{H_{zz}}H,\qquad
 \mathcal A_Hf=\frac1{2\pi i}\int_\Gamma Z_H(w)Q_f(w,z)\,dw.
\]
The divided difference $Q_f$ is a polynomial, including when $w=z$.

\begin{proposition}[Complete-source equation]\label{prop:source-coeff}
The prepared polynomial satisfies
\begin{equation}\label{eq:source-coeff}
 P_t=-\mathcal A_HP=-P_{zz}-2\mathcal T_P(bP_z),
 \qquad b=U_z/U,
\end{equation}
where
\[
 \mathcal T_Pg(z)=\frac1{2\pi i}\int_\Gamma
                   \frac{g(w)}{P(w)}Q_P(w,z)\,dw.
\]
These identities remain valid at internal multiple zeros.
\end{proposition}

\begin{proof}
The operator $\mathcal T_P$ gives the degree less than $N$ analytic remainder
modulo $P$; in particular it fixes every polynomial of degree less than $N$
and annihilates $P$ times a function holomorphic in the disk. These statements
follow directly from Cauchy's formula and do not require simple roots.
Since $H_t/H=P_t/P+U_t/U$, multiplying by $Q_P$ and integrating gives
$P_t=\int_\Gamma(H_t/H)Q_P/(2\pi i)$. The heat equation proves the first
identity. Finally
$H_{zz}/U=P_{zz}+2bP_z+(U_{zz}/U)P$, and analytic remainder gives the second.
\end{proof}

Put $\zeta=z-z_0$, and let $S$ be the plain coefficient shift:
$S\zeta^j=\zeta^{j-1}$ for $j\ge1$, and $S1=0$. It is not
$d/dz$. Define the complete source moments
\[
 M_k(t)=\frac1{2\pi i}\int_\Gamma(w-z_0)^k Z_H(t,w)\,dw.
\]
Expanding a divided monomial gives, on polynomials of degree at most $n$,
\begin{equation}\label{eq:source-shift}
 \mathcal A_H=\sum_{j=1}^{n}M_{j-1}S^j.
\end{equation}
Thus all these operators commute, including operators constructed from
different source fields. They lower degree and are nilpotent.

\begin{theorem}[Finite coefficient propagator]\label{thm:nilpotent-propagator}
Let $m_j(t,s)=\int_t^sM_{j-1}(u)\,du$. The backward propagator of
$E'=-\mathcal A_HE$ on degree at most $n$ is
\begin{equation}\label{eq:nilpotent-propagator}
 \mathcal U_H(t,s)=\exp\!\left(\sum_{j=1}^n m_j(t,s)S^j\right).
\end{equation}
The exponential terminates after degree $n$. For the fixed coefficient norm
$\|f\|_R=\sum_j|f_j|R^j$, its exact induced norm is
\begin{equation}\label{eq:propagator-norm}
 B_n(t,s)=\sum_{k=0}^n
 \left|[v^k]\exp\!\left(\sum_{j=1}^n\frac{m_j(t,s)}{R^j}v^j\right)\right|.
\end{equation}
If $V=\int_t^s\sup_\Gamma|Z_H(u,w)|\,du$, then
\begin{equation}\label{eq:propagator-majorant}
 B_n(t,s)\le\sum_{\ell=0}^n\binom n\ell\frac{V^\ell}{\ell!}.
\end{equation}
\end{theorem}

\begin{proof}
Equation~\eqref{eq:source-shift} makes the time generators commuting
polynomials in $S$, so differentiating the finite exponential proves
\eqref{eq:nilpotent-propagator}. If $\mathcal U=\sum b_kS^k$, the weighted
column sum of its $j$th column is $\sum_{k=0}^j|b_k|R^{-k}$, maximal at
$j=n$. This proves the exact norm formula. The wall estimate gives
$|m_j|/R^j\le V$. Majorize the generating series by
$\exp(Vv/(1-v))$ through degree $n$. Summing those coefficients and using
$\sum_{k=\ell}^n\binom{k-1}{\ell-1}=\binom n\ell$ proves
\eqref{eq:propagator-majorant}.
\end{proof}

The estimate is polynomial in acquired integrated source sizes at fixed
degree. It is not asserted to be polynomial uniformly in height, population
size, or a late-time parameter. For a difference of monic degree-$N$
polynomials, use $n=N-1$ in the error propagator, while propagating the
polynomial itself uses $n=N$.

\subsection{Reference errors and what must be acquired}

Let $G$ be a reference holomorphic field with the same zero count and with
$G(u,w)\ne0$ for every $u\in[t,T]$ and $w\in\Gamma$, with complete monic prepared polynomial $P_*$. It need not be an exact
heat solution. Define
\[
 d_*=\partial_tP_*+\mathcal A_GP_*,\qquad E=P-P_*,\qquad
 r=(\mathcal A_H-\mathcal A_G)P_*+d_*.
\]
Then $E'=-\mathcal A_HE-r$, and hence
\begin{equation}\label{eq:paid-duhamel}
 E(t)=\mathcal U_H(t,T)E(T)+\int_t^T\mathcal U_H(t,s)r(s)\,ds.
\end{equation}
The complete reference defect is also a wall quantity:
\[
 d_*(z)=\frac1{2\pi i}\int_\Gamma
                   \frac{G_t+G_{zz}}G(w)Q_{P_*}(w,z)\,dw.
\]
If $\eta_Z\ge\sup_\Gamma|Z_H-Z_G|$ and
$W_R(P_*)=\sum_jj|p_{*,j}|R^j$, a closed finite error estimate is
\begin{equation}\label{eq:paid-error}
 \|E(t)\|_R\le B_{N-1}(t,T)\|E(T)\|_R+
 \int_t^TB_{N-1}(t,s)
       \bigl(\eta_Z(s)W_R(P_*(s))+\|d_*(s)\|_R\bigr)\,ds.
\end{equation}
A directly evaluated correlated residual can replace this absolute majorant.
Its size must be acquired independently; an unknown coefficient difference
inside a forcing bound is not free data.

For example, acquired field and second-derivative errors $e_0,e_2$, reference
wall floor $m_*>e_0$, and reference second-derivative upper bound $M_{2,*}$ give
\begin{equation}\label{eq:source-quotient-budget}
 \eta_Z\le\frac{e_2}{m_*-e_0}
              +\frac{M_{2,*}e_0}{m_*(m_*-e_0)}.
\end{equation}
An analytic field-error bound on a spatial collar of width $d$ gives
$e_2\le2e_{\rm collar}/d^2$. This Cauchy step concerns the field error,
not a logarithmic ratio continued through different interior zero sets.

For exact heat fields the integrated moments have the endpoint representation
\begin{equation}\label{eq:source-endpoint-log}
 m_j(t,T)=-\frac1{2\pi i}\int_\Gamma(w-z_0)^{j-1}
              \log\frac{H(T,w)}{H(t,w)}\,dw,
\end{equation}
using the logarithm lifted continuously from time $t$. Thus equally informed
endpoint wall data already determine the homogeneous propagator. The
commutation theorem does not confer an intrinsic advantage on a curved path;
a gain would have to come from improved source estimates or conditioning along
that path, with its entire error paid.

\begin{remark}
The fixed-source generator is different from free polynomial heat evolution
on arbitrary error vectors. For $H=z^3-6tz$, one finds
$\mathcal A_H=6S^2$ on degree at most three, while
$\partial_z^2(z^2)=2$. Comparing two independently evolving polynomial heat
fields includes the source-quotient difference in
\eqref{eq:paid-duhamel}. Conversely, suppressing the external unit in
\eqref{eq:source-coeff} changes the complete problem. Note~12 contains the
free-heat comparison, complex-path formulation and collision controls.
\end{remark}

\section{A complete reference at microscopic real heat times}
\label{sec:microscopic-reference}

We now specify a theta source input rather than assume a small wall error.
Throughout,
\[
 H_t(z)=\int_0^\infty e^{tu^2}\Phi(u)\cos(zu)\,du,
 \qquad H_0(z)=\tfrac18\xi(\tfrac12+\tfrac{iz}2).
\]
The effective approximation of Polymath~\cite{polymath} applies for real
$0<t\le1/2$, $x\ge200$, $0\le y\le1$. In this section $z=x+iy$ and
\[
 N=\left\lfloor\sqrt{x/(4\pi)+t/16}\right\rfloor.
\]
We work on a chart where $N$ stays fixed. Changes of $N$ require explicit
patching; a jumping truncation is not a holomorphic reference.

Define, with principal logarithms away from their cuts,
\[
 M_0(s)=\frac{s(s-1)}{16}\pi^{-s/2}\sqrt{2\pi}
 \exp\!\left[(s/2-1/2)\operatorname{Log}(s/2)-s/2\right],
\quad
 \alpha(s)=\frac1{2s}+\frac1{s-1}+\frac12\operatorname{Log}\frac{s}{2\pi}.
\]
Here $\alpha=M_0'/M_0$. Put $s=(1-iz)/2$ and
\begin{equation}\label{eq:complete-finite-reference}
 G_{t,N}(z)=\sum_{\sigma\in\{s,1-s\}}\sum_{n=1}^N
 M_0(\sigma)n^{-\sigma}
             \exp\!\left[\frac t4(\alpha(\sigma)-\log n)^2\right],
 \qquad B_t(z)=M_0(s)e^{t\alpha(s)^2/4}.
\end{equation}
This rewrites the complete two-sum reference without complex conjugates.
In particular $G$ is locally holomorphic in $t,z$, although it is not an exact
heat solution.

\subsection{Explicit field and heat-defect budgets}

For completeness, a finite coarsening of the imported error bounds is as
follows. Set $X=x/(4\pi)$, $L=\log X$, and
\[
 p=\operatorname{Re}(s+t\alpha(s)/2),\quad
 \rho=e^{0.02y}X^{-y/2},\quad \nu=\frac{ty}{2(x-6)},
\]
\[
 S_0=\sum_{n=1}^N(1+\rho N^\nu n^y)e^{t\log^2n/4}n^{-p},
 \quad b_0=\frac{t^2L^2/16+0.626}{x-6.66},\quad
 E_{AB}=S_0(e^{b_0}-1).
\]
The exact value of $p$ avoids relying on slightly different printed
lower-bound corrections for its real part. Theorem~1.3 of
\cite{polymath} gives $|H_t-G_{t,N}|/|B_t|\le E_{AB}+E_{C,0}$, with
\[
 E_{C,0}=X^{-(1+y)/4}
 \exp\!\left[-\frac{tL^2}{16}
 +\frac{1.24(3^y+3^{-y})}{N-0.125}
 +\frac{3|L+i\pi/2|+10.44}{x-12}\right].
\]
For $tL^2\le K$ and $y=0$, $E_{AB}=O_K(X^{-3/4})$, whereas
$E_{C,0}=O_K(X^{-1/4})$. Indeed the retained weights have modulus at most
a constant times $n^{-1/2}$, and their relative error bracket is $O_K(1/x)$.
These are usable absolute errors even when $tL^2\to0$; the late-clock
assumption $x\ge e^{C/t}$ is not needed. That assumption is needed for the
much shorter two-carrier asymptotic, whose omitted absolute sum is not small
on the microscopic clock.

For a term $g=M_0(\sigma)n^{-\sigma}e^{t\beta^2/4}$, with
$\beta=\alpha(\sigma)-\log n$, direct differentiation gives
\begin{equation}\label{eq:finite-reference-defect}
 (\partial_t+\partial_z^2)g=-\frac g4
 \left[\alpha'+t\beta^2\alpha'
 +\frac t2\bigl((\alpha')^2+\beta\alpha''\bigr)
 +\frac{t^2}{4}\beta^2(\alpha')^2\right].
\end{equation}
Here both $\sigma=s$ and $\sigma=1-s$ satisfy
$\partial_z^2=-\partial_\sigma^2/4$. Taking absolute values in this finite
sum is an explicit defect budget. Since $\beta=O(\log x)$,
$\alpha'=O(x^{-1})$, $\alpha''=O(x^{-2})$, it proves
\begin{equation}\label{eq:uncorrected-defect-order}
 |G_t+G_{zz}|/|B_t|=O_K(X^{-3/4})\qquad(y=0,\ tL^2\le K).
\end{equation}
No source quotient has been divided by an unacquired wall floor in this step.

\subsection{An analytic continuation of the explicit correction}

Corollary~6.4 and Proposition~6.6(vi) of \cite{polymath} improve the field
remainder by retaining a correction $C_t$. Define
\[
 v_0=x/2+\pi t/8,\quad p_N(v)=1+2N-\sqrt{2v/\pi},
\]
\[
 U(v)=\exp\!\left[-i\left(\frac v2\log\frac v{2\pi}-\frac v2-\frac\pi8\right)\right],
 \quad
 C_0(p)=\frac{e^{\pi i(p^2/2+3/8)}-i\sqrt2\cos(\pi p/2)}{2\cos(\pi p)}.
\]
The singularities at $p=\pm1/2$ are removable, with value $(1-i)/4$.
The displayed source correction is
\[
 C_t(x+iy)=2(-1)^Ne^{-\pi iy/8+t\pi^2/64}
       \operatorname{Re}\{M_0(iv_0)C_0(p_N(v_0))U(v_0)e^{\pi i/8}\}.
\]
It gives
\begin{equation}\label{eq:corrected-field-remainder}
 \frac{|H_t-(G_{t,N}-C_t)|}{|B_t|}\le E_{AB}+E_C,
\end{equation}
where
\[
 E_C=X^{-(1+y)/4}
 e^{-tL^2/16+(3|L+i\pi/2|+3.58)/(x-8.52)}
 \left[\frac{1.24(3^y+3^{-y})}{N-0.125}+\frac{6.92}{x-12}\right].
\]
The extra $1/N$ makes this remainder $O_K(X^{-3/4})$ on the real axis.
However, the displayed $C_t$ is generally not holomorphic in $z$; its real-part
operation cannot be differentiated as an analytic reference.

There is an explicit analytic replacement. Let
$C_0^\#(p)=\overline{C_0(\bar p)}$,
$C_R=(C_0+C_0^\#)/2$, $C_I=(C_0-C_0^\#)/(2i)$, and
\[
 V_N(v)=v^{3/2}C_R(p_N(v))-v^{1/2}C_I(p_N(v)),\qquad
 k_N=-(-1)^N\frac{\sqrt\pi}{4}.
\]
The functions $C_R,C_I$ are holomorphic continuations of the real and imaginary
parts on the real axis, not real-part operations on complex arguments. Set
\begin{equation}\label{eq:analytic-correction}
 \widetilde C_{t,N}(z)=k_Ne^{-\pi z/8-\pi^2t/64}
                      V_N(z/2+\pi t/8),\qquad
 \widehat G=G-\widetilde C.
\end{equation}
Use principal square roots on $\operatorname{Re}v>0$.

\begin{proposition}[Corrected analytic reference]\label{prop:corrected-reference}
The correction \eqref{eq:analytic-correction} agrees with $C_t$ at $y=0$.
With
\[
 Q(t,z)=\frac{\sqrt\pi}{4}\frac{e^{-\pi x/8-\pi^2t/64}}{|B_t(z)|},
\]
it satisfies the explicit bounds
\begin{align}
 \frac{|H_t-\widehat G|}{|B_t|}
 &\le E_{AB}+E_C+Q(t,z)\frac{|y|}{2}
       \sup_{0\le\theta\le1}|V_N'(v_0+i\theta y/2)|,
                                                   \label{eq:analytic-correction-error}\\
 (\partial_t+\partial_z^2)\widetilde C
 &=\frac{k_N}{4}e^{-\pi z/8-\pi^2t/64}
                         V_N''(z/2+\pi t/8).       \label{eq:analytic-correction-defect}
\end{align}
For fixed $K$ and $y_0<1$, its normalized field allowance is
$O_{K,y_0}(X^{-(3+y)/4})$ uniformly for $0\le y\le y_0$ and
$0<tL^2\le K$. Its correction defect is
$O_{K,y_0}(X^{-(5+y)/4})$. In particular both the field error and total heat
defect of $\widehat G$ are $O_K(X^{-3/4})$ at $y=0$.
\end{proposition}

\begin{proof}
For positive real $v$, substituting
$\operatorname{Log}(iv/2)=\log(v/2)+i\pi/2$ gives the exact cancellation
\[
 M_0(iv)U(v)e^{\pi i/8}
       =-\frac{\sqrt\pi}{8}\sqrt v\,(v+i)e^{-\pi v/4}.
\]
Consequently
$C_t(x+iy)=k_Ne^{-\pi(x+iy)/8-\pi^2t/64}V_N(v_0)$.
The fundamental theorem of calculus from $v_0$ to $v_0+iy/2$ proves
\eqref{eq:analytic-correction-error}. In the heat defect, the coefficients of
$V_N$ and $V_N'$ cancel because the argument derivatives are $1/2$ and
$\pi/8$, while the exponential logarithmic derivatives are $-\pi/8$ and
$-\pi^2/64$. This proves \eqref{eq:analytic-correction-defect}.

For completeness the derivative sizes have explicit, zero-independent
constants. Along the comparison segment,
\[
 |p_N(v)-p_N(v_0)|\le\frac{|y|}{2\sqrt{2\pi v_0}}<\frac18,
\]
so $|p_N(v)|<9/8$. On this disk both $C_R^{(j)}$ and $C_I^{(j)}$ are bounded
by $K_j=j!8^jM_*$ for $j=0,1,2$, where
\[
 M_*=\frac{e^{25\pi/32}+\sqrt2\cosh(5\pi/8)}{2\sinh(\pi/8)}.
\]
Indeed $|\cos\pi p|\ge\sinh(\pi/8)$ on $|p|=5/4$: split into
$|\operatorname{Im}p|\ge1/8$ and its complement, where
$1<|\operatorname{Re}p|\le5/4$. Bound the numerator on that circle, remove
$\pm1/2$, and apply Cauchy's estimate with radius $1/8$.
Since $p_N'=-1/\sqrt{2\pi v}$ and $p_N''=-p_N'/(2v)$, differentiation gives
$V_N'=O(|v|)$ and $V_N''=O(|v|^{1/2})$ with these explicit constants.
Finally $Q=O_{K,y_0}(X^{-(7+y)/4})$ follows directly from the displayed
Stirling factor. Combine these estimates with
\eqref{eq:finite-reference-defect}. Note~14 records the full derivative
formulas and convenient finite majorants.
\end{proof}

\subsection{The remaining wall and time gates}

These are absolute source allowances. To use the coefficient theorem, acquire
a wall floor $m_f>e_0$ for $f=\widehat G/B$, where $e_0$ bounds the
right-hand side of \eqref{eq:analytic-correction-error}. Then
$|H/\widehat G-1|\le e_0/m_f$. For a common analytic normalization write
$e=H/B-f$ and $\lambda=B'/B$. The complete quotient cancellation is
\[
 \frac{H_{zz}}H-\frac{\widehat G_{zz}}{\widehat G}
 =\frac{fe''-ef''+2\lambda(fe'-ef')}{f(f+e)}.
\]
Thus an acquired spatial collar supplies derivative errors without introducing
internal root gaps. The local reference derivative ratios still require
bounds. A collar crossing the real axis must convert the reflected error to
one analytic carrier before applying Cauchy's estimate.

All imported field errors here concern real $t$. The explicit references
extend analytically in complex time on their charts, but that does not extend
the error for $H$ with the same bound. A complex-time route needs a separately
paid source remainder. Finally a fixed-$N$ chart is not the same as a fixed
number of roots in a selected block: a prepared population count, common wall,
and the forest denominator remain separate acquisition obligations. The finite
certificate in Note~13 pays those obligations in its original bounded box.
No estimate in this section supplies the missing unbounded-height small-gap
upper bound or a global simplicity conclusion.

\section{Complete-source acquisition on a fixed theta box}
\label{sec:finite}
The following calculation supplies actual canonical source data to the
coefficient interface. Its purpose is a bounded calibration, not a new
high-height zero result. The parameters were fixed before acquisition:
\begin{equation}
 \Omega=[26,46]+i[-2,2],\qquad
 D=[-10^{-4},10^{-4}]+i[-10^{-4},10^{-4}].
 \label{eq:finitebox}
\end{equation}
Every zero in the rectangle is included. No individual root location,
root disk, RH assumption or internal gap is an input.

\begin{theorem}[Finite complete-theta certificate]
\label{thm:finitecertificate}
For the unchanged field \eqref{eq:canonical} on \eqref{eq:finitebox},
the interval computation and analytic error bounds below establish
\[
 \inf_{D\times\partial\Omega}|H|>4.94\cdot10^{-7},\qquad
 \inf_{D\times\Omega}(|H|^2+|H_z|^2)>1.1178\cdot10^{-13}.
\]
The complete spatial population has multiplicity count two for every
\(\tau\in D\). Its centered monic quadratic has discriminant satisfying
\[
 189.68<\Re\Delta(\tau)<189.80,\qquad
 |\Im\Delta(\tau)|<0.058.
\]
Hence both roots are distinct throughout \(D\), and real and simple for
every real \(t\in[-10^{-4},10^{-4}]\). At \(t=10^{-4}\),
\[
 4.2149\cdot10^{-6}<q(t)<4.2176\cdot10^{-6}<\tfrac12.
\]
These conclusions concern only this finite population.
\end{theorem}

\subsection{Uniform error for the complete source}
Arb ball arithmetic \cite{arb}, through python-flint 0.9.0 at 192-bit
precision, encloses the first 96 Taylor coefficients of
\(H_0(z)=\xi(1/2+iz/2)/8\) about the real center 36. Real-type symmetry
justifies real coefficients; the computed imaginary coefficient balls
are checked to contain zero. Denote the degree-95 polynomial by
\(P_{95}\). Positivity of \(\Phi\) gives
\[
 |H_0(36+w)|\le H_0(24i)\quad(|w|\le24),
\]
so the full spatial Taylor remainder on the radius-12 disk is at most
\begin{equation}
 E_{\mathrm{sp}}=2H_0(24i)2^{-96}.
 \label{eq:spatialtail}
\end{equation}
On the radius-11 disk, Cauchy's inequality on a radius-one circle gives
the derivative error \(j!E_{\mathrm{sp}}\). The physical rectangle lies
strictly inside that disk since \(\sqrt{104}<11\).

Put \(T_*=1/5000\), so the entire time square, including its outward
interval representation, lies in \(|\tau|<T_*\). Use
\begin{equation}
 G_\tau(z)=\sum_{k=0}^{5}\frac{(-\tau)^k}{k!}P_{95}^{(2k)}(z).
 \label{eq:finiteheat}
\end{equation}
For \(|\Im z|\le2\), Taylor's integral remainder for the source
exponential and \eqref{eq:spatialtail} give
\begin{equation}
 |H_j(\tau,z)-G_\tau^{(j)}(z)|
 \le E_{\mathrm{sp}}\sum_{k=0}^{5}
       \frac{T_*^k(j+2k)!}{k!}
       +\frac{T_*^6}{6!}M_{12+j},\qquad 0\le j\le3,
 \label{eq:heaterror}
\end{equation}
where
\[
 M_j=\int_0^\infty u^je^{T_*u^2}\Phi(u)\cosh(2u)\,du.
\]
Every source index and the entire integration tail are paid. In fact,
with
\[
 C_\Phi=\frac{2\pi^2}{1-16e^{-3\pi}},\quad
 B_j=T_*/e^2+(11+j)/(2e),\quad X_*=(\pi/2)e^4,
\]
one has
\begin{equation}
 M_j\le e^{T_*}H_0(2i)+\frac{C_\Phi}{4X_*}
                   \exp\{B_j^2/(2\pi)-X_*\}.
 \label{eq:momentbound}
\end{equation}
For \(u\le1\), retain \(\cosh(2u)\) and bound \(u^je^{T_*u^2}\)
by \(e^{T_*}\). For \(u\ge1\), use
\(\Phi\le C_\Phi e^{9u-\pi e^{4u}}\), \(u^j\le e^{ju}\),
\(u^2\le e^{2u}/e^2\), and \(u\le e^{2u}/(2e)\).
Completing a square in \(e^{2u}\), followed by
\(\Gamma(0,X_*)\le e^{-X_*}/X_*\), proves \eqref{eq:momentbound}.
All four errors in \eqref{eq:heaterror} are below \(6\cdot10^{-27}\).
This is an approximation to the complete field, not a replacement source.

\subsection{Wall count and coefficient enclosure}
On each accepted wall segment \(\gamma_\nu\), the source oracle gives
\[
 m_\nu\le\inf|H_0|,\qquad
 \omega_\nu\ge\sup_{D\times\gamma_\nu}|H_\tau-H_0|,
 \qquad \eta_\nu=\omega_\nu/m_\nu<1/4.
\]
The time difference is evaluated with its constant heat term cancelled
before interval evaluation. Sixteen segments, found in 28 cell calls,
cover the complete closed rectangle. The largest \(\eta_\nu\) is less
than \(3.82\cdot10^{-4}\). Thus Rouch\'e's theorem transfers the
entire multiplicity count from the anchor.

At the anchor, validated integration encloses
\begin{equation}
 p_k(0)=\frac1{2\pi i}\oint_{\partial\Omega}
       (w-36)^k\frac{H_0'(w)}{H_0(w)}\,dw,\qquad k=0,1,2.
 \label{eq:anchormoments}
\end{equation}
The meromorphic quotient formed from the entire finite polynomial is passed
to the validated integrator. It is holomorphic near the certified wall;
the integrator checks its required complex neighborhoods. A physical-cell range enclosure, which is valid only on its
declared physical subset, is not an analytic callback.
If \(|H_0|\ge m\), \(|H_0'|\le L\), and source errors are
\(e_0,e_1\), the external correction is
\begin{equation}
 \left|\frac{H_0'}{H_0}-\frac{G_0'}{G_0}\right|
 \le\frac{e_1}{m}+\frac{(L+e_1)e_0}{m(m-e_0)},\qquad m>e_0.
 \label{eq:quotienterror}
\end{equation}
Here \(e_0=e_1=E_{\mathrm{sp}}\). Integrating this bound on the
accepted segments pays the complete-source correction.
The resulting \(p_0\) enclosure isolates the unique integer two.
The argument principle therefore supplies the count rather than a
numerical guess about root locations.

The principal relative logarithm
\(\ell_\tau=\operatorname{Log}(H_\tau/H_0)\) is single-valued around
the wall and obeys \(|\ell_\tau|\le-\log(1-\eta_\nu)\).
Integration by parts gives
\begin{align}
 p_k(\tau)-p_k(0)
   &=-\frac{k}{2\pi i}\oint(w-36)^{k-1}\ell_\tau(w)\,dw,\nonumber\\
 |p_k(\tau)-p_k(0)|
   &\le\frac{k}{2\pi}\sum_\nu
      |\gamma_\nu|r_\nu^{k-1}[-\log(1-\eta_\nu)],
 \label{eq:logtransfer}
\end{align}
where \(r_\nu=\sup_{\gamma_\nu}|w-36|\). No internal gap or
interior logarithmic ratio is differentiated. The acquired values are
\[
 p_1(0)=-1.6864704389875024\pm7.4\cdot10^{-12},\quad
 p_2(0)=96.29229323297757\pm4.1\cdot10^{-11},
\]
with these bounds on each complex component. The uniform continuation
errors are below \(0.001772\) and \(0.025736\), respectively.
Since \(\Delta=2p_2-p_1^2\), outward interval arithmetic proves the
discriminant enclosure in Theorem~\ref{thm:finitecertificate}.
Conjugation symmetry gives real coefficients on real time slices;
positivity of the discriminant then supplies reality and simplicity.
The quadratic forest formula gives \(q=8t/(\Delta+8t)\).

For the first-jet floor, 31 closed spatial rectangles, obtained in
61 cell calls, cover \(\Omega\) while every cell uses all of \(D\).
The sum of the squared lower bounds on \(|H|\) and \(|H_z|\) is
positive on every cell. Exact rational area, containment and disjointness
checks verify the cover. This is also a direct quantitative certificate
of regularity on every physical cap and face.

\subsection{Matched direct comparison and reproducibility}
Direct evaluation of \eqref{eq:anchormoments} at \(t=10^{-4}\), using
the same source, wall, tolerance and correction, gives the much tighter
ball
\[
 q(10^{-4})=
 4.2162799666292433\cdot10^{-6}\ \pm\ 2.34\cdot10^{-18}.
\]
The uniform complex-time transport ball has radius
\(1.28\cdot10^{-9}\). The different information delivered by the two
calculations should be kept in view: the latter controls all of \(D\),
whereas the former is a direct terminal enclosure. For real-time
simplicity alone, four source values give sign changes in \((28,29)\)
and \((42,43)\), uniformly on the real interval; together with the
shared complete count they give a cheaper direct conclusion.
No practical advantage for continuation is established here.

Research Note~13 and Review~11 contain the detailed proof and audit.
The sibling \texttt{numerics} folder contains the source core,
the \texttt{07} source and wall scripts and their small JSON receipts,
and the \texttt{08} complete first-jet cover. The installed wall replay
passes 109 assertions, the independent complete-source quadrature
checks pass 46, and the full-volume cover passes five. The source checks
include four direct theta-integral evaluations with both omitted modes
and the integration tail paid. They supplement the uniform analytic
bounds, rather than proving them by sampling.
The receipts record numerical precision, source errors, geometry,
script hashes and timings. Typical wall work is a fraction of a second;
the anchor and direct contour integrations take several seconds on the
recorded machine. These timings are implementation observations, not
asymptotic complexity claims.

\section{Complete physical contours and matched comparisons}
\label{sec:complete-contours}

A contour calculation must preserve the observable before it estimates its
pieces.  Write $\tau=t+i\sigma$, $z=x+iy$, and $H_j=\partial_z^jH$ for a
holomorphic solution of $H_\tau=-H_{zz}$.  For the first jet
$G=(H_0,H_1)$, put
\[
 E=|H_0|^2+|H_1|^2,\qquad
 L=H_1^2-H_0H_2,\quad N=H_2^2-H_1H_3,\quad M=H_1H_2-H_0H_3.
\]
Where $E>0$, the normalized sphere-volume form is
\[
 \mathcal C=\frac{\mathcal A\wedge d\mathcal A}{4\pi^2},\qquad
 \mathcal A=\frac{\operatorname{Im}(G^*dG)}{E}.
\]
Define
\[
 (f,p)=\frac{(\Re(L\overline N),\Im(L\overline N))}{2\pi^2E^2},
 \qquad
 (r,s)=\frac{(\Re(M\overline N),\Im(M\overline N))}{2\pi^2E^2}.
\]
Direct evaluation on the coordinate vectors, using
$G_t=(-H_2,-H_3)$, $G_\sigma=iG_t$, $G_x=(H_1,H_2)$ and
$G_y=iG_x$, gives
\begin{equation}
 \mathcal C=f\,d\sigma\wedge dx\wedge dy
 -p\,dt\wedge dx\wedge dy
 +r\,dt\wedge d\sigma\wedge dy
 -s\,dt\wedge d\sigma\wedge dx.
 \label{eq:physical-density}
\end{equation}
These formulas use no projective denominator such as $H$ or $H_z$ alone.
A zero of either coordinate is harmless when $E>0$.

Let $\Omega=[A,B]+i[-b,b]$ and
$W_-=[-\epsilon,0]\times[-Y,Y]\times\Omega$.  Give the cap at fixed $t$
the forward orientation $d\sigma\wedge dx\wedge dy$, and write
\[
 I_t=\int_{-Y}^Y\int_A^B\int_{-b}^b f(t,\sigma,x,y)\,dy\,dx\,d\sigma.
\]
The entire remaining outward-face contribution is
\begin{align}
 S_\epsilon={}&\int_{-\epsilon}^0\int_A^B\int_{-b}^b
 [p(t,Y,x,y)-p(t,-Y,x,y)]\,dy\,dx\,dt\notag\\
 &+\int_{-\epsilon}^0\int_{-Y}^Y\int_{-b}^b
 [r(t,\sigma,B,y)-r(t,\sigma,A,y)]\,dy\,d\sigma\,dt\notag\\
 &+\int_{-\epsilon}^0\int_{-Y}^Y\int_A^B
 [s(t,\sigma,x,b)-s(t,\sigma,x,-b)]\,dx\,d\sigma\,dt.
 \label{eq:all-physical-faces}
\end{align}

\begin{proposition}[Complete physical-cap identity]
Suppose the first jet is nonzero on $\partial W_-$ and its interior joint
zeros are isolated.  Then
\begin{equation}
 I_0=I_{-\epsilon}-S_\epsilon+Q_-,\qquad
 Q_-=\sum_{p\in\operatorname{int}W_-}\binom{m_p}{2}\ge0,
 \label{eq:complete-cap-transfer}
\end{equation}
where $m_p$ is the spatial multiplicity of the heat zero at $p$.
In particular $I_{-\epsilon}-S_\epsilon$ is a lower bound for the same
whole cap $I_0$.
\end{proposition}

\begin{proof}
Remove small four-dimensional balls about the joint zeros.  On the
remaining domain $\mathcal C$ is the pullback of a closed three-form on
$S^3$.  Stokes' theorem therefore identifies the complete outer boundary
integral with the sum of the local degrees on small spheres oriented outward
from the removed balls, opposite to their induced orientations as holes.  The
local holomorphic heat-degree calculation assigns a spatial multiplicity
$m$ the positive degree $\binom m2$; this is the local result used in
Investigation Note~02, Section~2.  Substituting the outward orientations in
\eqref{eq:physical-density} yields \eqref{eq:all-physical-faces} and
\eqref{eq:complete-cap-transfer}.
\end{proof}

A joint zero on a cap does not meet this proposition's hypotheses.
Removing it requires retaining the actual partial sphere and deleted cap
piece; no automatic half-integer or one-sided replacement is available.
Likewise, shared computational interfaces cancel only when they are the
same physical face with opposite orientation.  A changed time, radius or
population creates a joining contribution.  Optional root charts must be
supplemented by their full spatial complement.  Even a prepared
factorization $H=UP$ does not identify its discriminant with the physical
root product: for $N_0$ interior roots,
\[
 \prod_iH_z(\tau,z_i)
 =(-1)^{N_0(N_0-1)/2}\operatorname{disc}(P)\prod_iU(\tau,z_i).
\]
The nonzero unit remains part of the physical observable.

\subsection{Common time and a height graph are different constructions}

Along a common complex-time path $\tau=\gamma(s)$, a zero-free spatial wall
preserves the complete population, and its prepared coefficients remain
holomorphic through internal collisions.  Moving the wall does not change
its contour moments while the same complete population remains enclosed.
This fact does not control the spectrum of a pair matrix on a complex
slice.  If $R=(I+aL)^{-1}$ and $q=\Tr(I-R)$, then
\[
 \frac{dq}{ds}=\Tr\!\left[R(a'L+a\gamma'L_\tau)R\right].
\]
A matrix-norm estimate requires a lower singular-value bound for $I+aL$;
its two derivative terms should remain coupled.  Distinct roots alone do
not guarantee this resolvent condition.  For example,
$H=z^2-1-2\tau$ has $q=2\tau/(1+4\tau)$, whose pole at $-1/4$ is not a
root collision.  The coefficient/forest method elsewhere in this paper
provides a different scalar denominator test, with its own stated scope.

A height graph instead samples
$F(z)=H(h(\Re z)+i\sigma,z)$, with real $C^1$ function $h$.  It is generally
nonholomorphic:
\begin{equation}
 F_z=H_z-\tfrac12h'H_{zz},\qquad
 F_{\bar z}=-\tfrac12h'H_{zz},\qquad
 \frac{dF}{F}=\frac{H_z}{H}\,dz-h'\frac{H_{zz}}H\,dx.
 \label{eq:height-graph-corrections}
\end{equation}
The resulting winding is a topological degree, rather than an unsigned
holomorphic zero count.  At a graph zero with $H_z\ne0$, its Jacobian is
\[
 J_F=|H_z|^2\Re\!\left(1-h'\frac{H_{zz}}{H_z}\right).
\]
Positive orientation therefore needs a separate transversality condition.
For isolated transverse graph zeros, weighted counting also has the bulk
correction
\[
 \sum_{F(\zeta)=0}\operatorname{ind}_\zeta(F)\zeta^k
 =\frac1{2\pi i}\oint_{\partial\Omega}z^k\frac{dF}{F}
 +\frac{k}{\pi}\int_\Omega z^{k-1}\frac{F_{\bar z}}F\,dA,
 \qquad k\ge1.
\]
This follows by Stokes' theorem on the punctured domain; the integrand is
locally integrable at a transverse zero.  A negative graph $t=-\epsilon(x)$
also changes the physical cap density to $f+\epsilon'(x)r$.
Neither correction is optional.  These distinctions and finite graph
comparison bounds are proved in Investigation Note~08, Sections~6--7.

\subsection{What the comparison has established}

The direct estimator receives the same complete source representation,
precision, signed correlations, wall data and any acquired factorization as
the deformed estimator.  A common analytic endpoint identity does not by
itself make a detour cheaper.  Under a uniform residual allowance $e$, a
path pays its length times $e$; a longer path can improve that budget only
if its actual source or conditioning bounds improve enough.  This is a
scoped comparison, not an exclusion of useful contour deformation.

The finite canonical test uses the original box
\[
 [26,46]+i[-2,2],\qquad
 \tau\in[-10^{-4},10^{-4}]+i[-10^{-4},10^{-4}].
\]
An interval certificate supplies $E>1.1178\cdot10^{-13}$ throughout the
whole box, with no deleted root disks.  Thus $Q_-=0$ and both estimators
know the exact equality in \eqref{eq:complete-cap-transfer}.  A single
bounded range-quadrature pilot, using 512 spatial leaves and all the faces
in \eqref{eq:all-physical-faces}, nevertheless gives only
\[
 I_0=2.0247388041\cdot10^{-5}\ \pm1.42\cdot10^{-4},\qquad
 I_{-\epsilon}-S_\epsilon
 =2.0247432516\cdot10^{-5}\ \pm1.42\cdot10^{-4}.
\]
These are certified enclosing intervals, not accurate midpoint estimates.
Both lower endpoints are negative.  The raw deformed width is slightly
larger, approximately $0.000283934$ versus $0.000283522$, and it requires
additional face evaluations.  The close $\sigma$ faces were subtracted
algebraically before interval bounds, retaining their factor $2iY$.
The total side-face interval has radius below $2.05\cdot10^{-7}$;
spatial range overestimation of the caps dominates the uncertainty.
Investigation Note~15 and its replay give the complete budgets and scope.

Accordingly this finite experiment pays the source, regularity and full
boundary-accounting obligations, but demonstrates neither a positive
complete cap lower bound nor a contour advantage.  It does not justify a
larger-height campaign.  A further numerical comparison, if undertaken,
should test higher-order correlated spatial integration on the same
observable and at a fixed total budget, with the same improvement made
available to the direct baseline.

% Provisional research module for main.tex, not book.tex.
% Prepared for Edward Baker, 23 October 2026.
% Drafted with OpenAI Codex, gpt-6-astra, reasoning effort ultra.
% Source: Research Note 16 and Review 07; no new canonical estimate is claimed.

\section{Provisional conservation laws for spatial contours}
\label{sec:conserved-contours}

This section records an additional research direction, rather than a
completed estimate for the canonical field.  The conservation identities,
stationary actions and finite-flow bounds below are exact.  Their useful
size at growing height, and any consequence for endpoint simplicity,
remain unproved.  In particular, no global graph, deformation or
collision-exclusion theorem for theta is asserted.

Here time is real, and the proposed contour is spatial:
\begin{equation}
 \Gamma_t:\quad z=T+i\epsilon(T,t),\qquad H_t=-H_{zz}.
 \label{eq:cc-geometry}
\end{equation}
The height coordinate $T$ and heat time $t$ are independent.  This is
different from a boundary at time $-\epsilon(T)$.  On a compact time
interval, assume that $H$ is holomorphic on a neighborhood of every
spatial region used.  When reflection or the canonical source is needed,
it is stated explicitly.  No curvature smoothing, restoring force or
random source is added.

\subsection{The heat current, primitive and positive metric}

Write a complex velocity as $u=u_1+iu_2$.  A material point on
\eqref{eq:cc-geometry} satisfies $\dot T=u_1$; the chain rule gives
\begin{equation}
 \epsilon_t+u_1\epsilon_T=u_2.
 \label{eq:cc-graph}
\end{equation}
The field itself supplies a particularly simple conservation law:
$H_t+\partial_zH_z=0$.

\begin{proposition}[Material preservation of the heat form]
\label{prop:cc-heat-form}
On a region where $H\ne0$, put $p=H_z/H$ and let
$Z_t=p(t,Z)$.  For every regular material arc,
\begin{equation}
 H(t,Z(a,t))Z_a(a,t)
 =H(t_0,Z(a,t_0))Z_a(a,t_0).
 \label{eq:cc-material-form}
\end{equation}
Thus every material subarc preserves $\int H\,dz$, and its graph obeys
\eqref{eq:cc-graph} with $u=p$.
\end{proposition}
\begin{proof}
The time derivative of the left side is
$[H_t+\partial_z(Hp)](t,Z)Z_a=0$.
More generally, differentiating any moving open-arc integral gives
\begin{equation}
 \frac{d}{dt}\int_{\Gamma_t}H\,dz
 =[H\dot z-H_z]_{\rm endpoints}.
 \label{eq:cc-endpoints}
\end{equation}
Consequently preservation of a single total integral constrains only
its endpoints.  The closed integral of entire $H$ is identically zero.
Local material preservation is the stronger requirement in
\eqref{eq:cc-material-form}.  The freedom $Hu=H_z+C(t)$ also preserves
that form; choosing the literal heat current fixes $C=0$.
\end{proof}

For an arbitrary moving graph, put
$g(T,t)=H(t,T+i\epsilon)(1+i\epsilon_T)$.  Direct differentiation gives
\begin{equation}
 g_t+\partial_T\{H_z(t,T+i\epsilon)-i\epsilon_tH(t,T+i\epsilon)\}=0.
 \label{eq:cc-graph-density}
\end{equation}
Under $u=p$, the expression in braces is $(\Re p)g$.
Integrals between material height endpoints are conserved; fixed height
endpoints retain this flux.

There is also a positive planar conservation law:
\begin{equation}
 \partial_t|H|^2+\operatorname{div}_{x,y}(H_z\overline H)=0,
 \qquad x=\Re z,\quad y=\Im z.
 \label{eq:cc-positive}
\end{equation}
Indeed the divergence of the displayed complex current, regarded as a
real vector, is $2\Re(H_{zz}\overline H)$.  Its velocity divided by
density is again $p$.  Reynolds transport therefore preserves
$\int_{D_t}|H|^2\,dx\,dy$ for compact material domains with regular
swept boundaries.  Only a boundary tube needs a regular velocity;
interior zeros need not be excluded.  Taking absolute values in
\eqref{eq:cc-material-form} also preserves weighted length
$\int_{\Gamma_t}|H|\,|dz|$ and the metric $|H|^2|dz|^2$.
These weights degenerate at zeros: an ordinary length bound still needs
a positive lower bound for $|H|$ along the boundary.

For even $H$, define the odd primitive, distinct from the prepared
polynomial $P$ elsewhere in this manuscript, by
\begin{equation}
 \mathcal P(t,z)=\int_0^zH(t,w)\,dw,
 \qquad \mathcal P_z=H,\qquad \mathcal P_t=-H_z.
 \label{eq:cc-primitive}
\end{equation}
The last identity follows from $H_z(t,0)=0$.  For a non-even field the
same identity is obtained by choosing a time-dependent additive
constant.  The material law is equivalently
\begin{equation}
 \mathcal P(t,Z(a,t))=\mathcal P(t_0,Z(a,t_0)).
 \label{eq:cc-primitive-image}
\end{equation}
It continues the inverse image of a prescribed target curve.  A local
inverse requires $H\ne0$; a global continuation also needs branch and
escape control.  In fact the canonical primitive has the finite
real-direction asymptotic value
\[
 \mathcal P(t,+\infty)=\frac12\int_{\R}H(t,x)\,dx
                     =\frac{\pi}{2}\Phi(0),
\]
independent of $t$, by \eqref{eq:cc-fourier} below.
Avoidance of critical values alone is therefore not a global inverse
theorem: a finite-buffer boundary-image separation or another
properness/escape estimate is needed.  Even a regular material curve
can cease to be a graph.  For example, a patch of
$\Im\mathcal P=c$ has
\[
 \epsilon_T=-\frac{\Im H}{\Re H},\qquad
 \epsilon_t=\frac{\Im H_z}{\Re H},
\]
which requires the stronger condition $\Re H\ne0$.

\subsection{Two Noether formulations and the adjoint hierarchy}

The primitive gives a geometric first-order action.  Write
$\mathcal P=X+iY$ and set
\begin{equation}
 \mathcal S_{\mathcal P}[Z]
 =\frac12\int\Im(\overline{\mathcal P(t,Z)}\,d\mathcal P(t,Z))
 =\frac12\int(X\,dY-Y\,dX).
 \label{eq:cc-primitive-action}
\end{equation}
The curvature of its spacetime one-form is
\[
 |H|^2dx\wedge dy+
 \Im(H_z\overline H)dt\wedge dx-
 \Re(H_z\overline H)dt\wedge dy.
\]
Its kernel with unit time component gives $\dot z=p$ when $H\ne0$.
Equivalently the target equations are $\dot X=\dot Y=0$.
Translations of $X$ and $Y$ change the Lagrangian by
$\dot Y/2$ and $-\dot X/2$, respectively.  Subtracting these endpoint
terms in Noether's formula gives charges $-Y$ and $X$.
This is a stationary action, linear in velocity; it supplies no positive
energy-minimization principle.

There is a complementary field-adjoint formulation
\cite{noether,ibragimov2007,anco}.  On a real spatial interval take
independent fields $H,\psi$ and
\begin{equation}
 \mathcal S[H,\psi]=\int(\psi H_t-\psi_xH_x)\,dx\,dt.
 \label{eq:cc-adjoint-action}
\end{equation}
Compactly supported variations give $H_t=-H_{xx}$ and
$\psi_t=\psi_{xx}$.  For complex fields, real and imaginary variations
of the real part of this action give the same equations.  The exact
scaling symmetry $(H,\psi)\mapsto(e^aH,e^{-a}\psi)$ has density
$\psi H$ and current $\psi H_x-\psi_xH$; explicitly,
\begin{equation}
 \partial_t(\psi H)+\partial_x(\psi H_x-\psi_xH)
 =\psi(H_t+H_{xx})+H(\psi_t-\psi_{xx}).
 \label{eq:cc-adjoint-identity}
\end{equation}
Thus the variational statement is proved here for the doubled system,
rather than assumed for a single-field conservative heat action.
For holomorphic adjoints it continues with $z$ derivatives.  Its
zero-additive-current material velocity is
\begin{equation}
 u_\psi=\frac{\psi H_z-\psi_zH}{\psi H}
 =p-\frac{\psi_z}{\psi},
 \qquad \psi H\ne0.
 \label{eq:cc-adjoint-velocity}
\end{equation}
The choice $\psi=1$ recovers the primary flow.  Polynomial adjoints are
$e^{t\partial_z^2}z^n$, beginning with
$1,z,z^2+2t,z^3+6tz,z^4+12tz^2+12t^2$.

\begin{proposition}[Canonical whole-line invariants]
\label{prop:cc-hierarchy}
Extend the source $\Phi$ in \eqref{eq:canonical} evenly to $\R$.
With $\widehat H_t(\xi)=\int_\R e^{-i\xi x}H(t,x)\,dx$, one has
\begin{align}
 e^{-t\xi^2}\widehat H_t(\xi)&=\pi\Phi(\xi),
 &&\xi\in\R,\label{eq:cc-fourier}\\
 I_n(t):=\int_\R(e^{t\partial_x^2}x^n)H(t,x)\,dx
 &=\pi i^n\Phi^{(n)}(0),
 &&n\ge0.\label{eq:cc-polynomial-invariants}
\end{align}
These integrations and their time derivatives are justified uniformly
on compact real time intervals.
\end{proposition}
\begin{proof}
The theta transformation makes the even extension of $\Phi$ smooth.
More explicitly, put
$\vartheta(v)=\sum_{m\in\mathbb Z}e^{-\pi m^2v}$ and
$h(u)=e^u\vartheta(e^{4u})$.
The identity $\vartheta(v)=v^{-1/2}\vartheta(1/v)$ makes $h$ even,
and direct differentiation of \eqref{eq:canonical} gives
$16\Phi=h''-h$.  Its positive-half-line expansion and all derivatives
decay superexponentially.  Consequently $e^{tu^2}\Phi(u)$ is Schwartz,
uniformly for $t$ in a compact interval, as is its Fourier transform.
The cosine normalization gives
$\widehat H_t=\pi e^{t\xi^2}\Phi$.
Finally $\int x^mH\,dx=i^m\widehat H_t^{(m)}(0)$, and the finite
expansion of $e^{t\partial_x^2}x^n$ gives
\[
 I_n=i^n\left.\partial_\xi^n
       (e^{-t\xi^2}\widehat H_t(\xi))\right|_{\xi=0}.
\]
This proves both formulas and pays all real-line endpoint terms.
\end{proof}

The continuum \eqref{eq:cc-fourier} determines $H$ by Fourier inversion.
Completeness of polynomial moments for arbitrary Schwartz data is not
asserted.  Both families have explicitly time-dependent weights:
they encode the known linear evolution and its uniqueness.  They do
not establish an autonomous Hamiltonian hierarchy, a Lax isospectral
structure or a coercive conserved energy for the contour.  The separate
KdV-hierarchy discussion in \cite{ibragimov2006} is not a theorem
imported into this problem.

\begin{proposition}[No common local material hierarchy]
\label{prop:cc-no-common-velocity}
On a nonempty open region with $H\ne0$, no single regular velocity
preserves every material subarc of all three forms
$H\,dz$, $zH\,dz$ and $(z^2+2t)H\,dz$.
\end{proposition}
\begin{proof}
Divide the last two material pullback identities by the first.
For each weight this forces $\psi_t+u\psi_z=0$.
The weight $z$ forces $u=0$, while $z^2+2t$ then forces $2=0$.
This argument allows arbitrary spatially constant current gauges.
\end{proof}
There is no conflict with simultaneous conservation of whole-line
integrals: that weaker statement has vanishing endpoint currents.
The hierarchy constrains the field supplying the velocity, rather than
giving infinitely many independent local boundary conditions.

Finally, on $H\ne0$, logarithmic differentiation gives the
Cole--Hopf relation \cite{cole,hopf}
\begin{equation}
 p_t+\partial_z(p_z+p^2)=0,
 \qquad U=2p:\quad U_t+UU_z=-U_{zz}.
 \label{eq:cc-burgers}
\end{equation}
The factor $2p$ in Burgers advection is not the $H\,dz$ material
velocity.  Preserving $p\,dz$ instead, with its literal current,
selects $(p_z+p^2)/p=H_{zz}/H_z$ where defined.  This transports the
value of $H$ and agrees at simple zeros with their usual heat velocity;
the sign is opposite to the forward-heat convention in
\cite{papanicolaou}.

\subsection{Regularity and a positive-source collision control}

The stationary solution $H=z$ has $p=1/z$ and
$Z(t)^2=Z(t_0)^2+2(t-t_0)$.  A point initially at $iy_0$ reaches a
simple zero after time $y_0^2/2$.  Thus the conserved material form does
not ensure a regular flow even in the absence of multiple zeros.

The following decaying example tests the hierarchy with positive source.
For $a,b>0$ and $d=1-4at>0$, set
\begin{equation}
 F(t,z)=d^{-1/2}e^{-az^2/d}
       \{1+e^{tb^2/d}\cos(bz/d)\}.
 \label{eq:cc-gaussian-control}
\end{equation}
The functions
$d^{-1/2}\exp(-az^2/d\mathbin{\pm}ibz/d+tb^2/d)$ solve backward heat
by differentiation, as does the unmodulated Gaussian.
Their sum proves $F_t=-F_{zz}$.
The initial Fourier transform is strictly positive:
\[
 \widehat F_0(\xi)=\sqrt{\pi/a}\left\{
 e^{-\xi^2/(4a)}+\tfrac12e^{-(\xi-b)^2/(4a)}
                  +\tfrac12e^{-(\xi+b)^2/(4a)}\right\}.
\]
This even, real-type, spatially entire family is uniformly Schwartz on
compact time intervals below $1/(4a)$; all its Fourier and polynomial
invariants are conserved.  For $0<t<1/(4a)$, its zero equation is
$\cos(bz/d)=-e^{-tb^2/d}\in(-1,0)$, so every zero is real and simple.
Indeed a nonreal $w$ with real $\cos w$ has either $\Im w=0$, or
$|\cos w|\ge1$.  At $t=0$ each $z=(2k+1)\pi/b$ is a double zero.
Therefore positive even Schwartz source, the complete hierarchy and
an interval of positive-time reality do not exclude endpoint doubles.
This family has different conserved values from theta and is defined
only below its indicated time singularity.  It tests a generic
implication; it does not exclude an estimate using the unchanged
canonical source and its actual constants.

\subsection{A regularized zero-current comparison}

A different conserved object gives a smooth finite-parameter model.
Fix $\rho,\lambda>0$ and define
\begin{equation}
 k_\rho(w)=\frac{\rho^2}{\pi(|w|^2+\rho^2)^2},\quad
 \mu_\rho=k_\rho(H)|H_z|^2,\quad
 j_\rho=k_\rho(H)H_{zz}\overline{H_z}=j_1+ij_2.
 \label{eq:cc-scalar-current}
\end{equation}
Pulling back $k_\rho(w)\,d\Re w\wedge d\Im w$ by $H(t,z)$ gives
\begin{equation}
 \Omega_\rho=\mu_\rho dx\wedge dy+j_2dt\wedge dx-j_1dt\wedge dy.
 \label{eq:cc-scalar-curvature}
\end{equation}
The target two-form is closed, so $d\Omega_\rho=0$, equivalently
$\partial_t\mu_\rho+\partial_xj_1+\partial_yj_2=0$.
No division by $H_z$ is used.
The outward upper-wall current per unit horizontal height is
\begin{equation}
 J_\rho=j_2-j_1\epsilon_T-\mu_\rho\epsilon_t.
 \label{eq:cc-relative-current}
\end{equation}
Choosing $\lambda\epsilon_t=J_\rho$ gives
\begin{equation}
 (\lambda+\mu_\rho)\epsilon_t+j_1\epsilon_T=j_2,
 \qquad u_\rho=\frac{j_\rho}{\lambda+\mu_\rho}.
 \label{eq:cc-compliant-flow}
\end{equation}
This transports the augmented density $\lambda+\mu_\rho$.
At fixed positive parameters the velocity is smooth on compact analytic
regions, including multiple zeros.  The parameter $\lambda$ is chosen,
not forced by heat; making it large suppresses motion by construction.

The same law has the exact spacetime path action $\int\mathcal A_\rho$,
where
\begin{equation}
 \mathcal A_\rho=\frac\lambda2(x\,dy-y\,dx)
 +\frac1{2\pi}\frac{\Im(\overline H\,dH)}{|H|^2+\rho^2}.
 \label{eq:cc-scalar-action}
\end{equation}
Differentiation gives
$d\mathcal A_\rho=\lambda dx\wedge dy+\Omega_\rho$.
For $V=\partial_t+u_1\partial_x+u_2\partial_y$, fixed-endpoint
variations require $\iota_Vd\mathcal A_\rho=0$, precisely
\eqref{eq:cc-compliant-flow}.  Cartan's formula then gives
$\mathcal L_V\mathcal A_\rho=d(\mathcal A_\rho(V))$.
Consequently complete closed material loops preserve circulation;
Stokes identifies it with augmented enclosed mass.  The full $dt$
component of the one-form is essential to this calculation.

Assume now real type, $\epsilon\ge0$, and a fixed height band $[A,B]$.
Then $j_2(t,T,0)=0$.  Define
\begin{align*}
 m(T,t)&=\int_0^{\epsilon(T,t)}\mu_\rho(t,T,y)\,dy,&
 q(T,t)&=\int_0^{\epsilon(T,t)}j_1(t,T,y)\,dy,\\
 M_\rho(t)&=\int_A^B m(T,t)\,dT,&
 E_{\rm end}(t)&=q(B,t)-q(A,t).
\end{align*}
The continuity equation and moving upper limit give
$m_t+q_T=-J_\rho$.  Hence, writing
$\bar\epsilon=(B-A)^{-1}\int_A^B\epsilon\,dT$,
\begin{equation}
 \partial_t(m+\lambda\epsilon)+\partial_Tq=0,
 \qquad
 \frac d{dt}\{M_\rho+\lambda(B-A)\bar\epsilon\}=-E_{\rm end}.
 \label{eq:cc-mean-balance}
\end{equation}
This retains the actual end currents; the canonical field is not
assumed periodic.  Positivity alone implies, for $t\ge t_0$,
\begin{equation}
 \bar\epsilon(t)\le\bar\epsilon(t_0)+
 \frac{M_\rho(t_0)+\int_{t_0}^t[-E_{\rm end}(s)]_+\,ds}
      {\lambda(B-A)}.
 \label{eq:cc-mean-bound}
\end{equation}
If independently $0\le M_\rho\le M_*$, the absolute change in mean is
at most $(M_*+|\int_{t_0}^tE_{\rm end}|)/(\lambda(B-A))$.
Neither statement controls roughness or repeated inward crossings:
\begin{equation}
 \int\!\int[-J_\rho]_+\,dT\,dt
 =\lambda\int\!\int[-\epsilon_t]_+\,dT\,dt.
 \label{eq:cc-one-sided-variation}
\end{equation}
Small signed mean change does not bound this positive variation.
The earlier nonreal population, height ends and endpoint coverage
remain necessary in any transfer to a multiplicity consumer.

\subsection{Conditional graph persistence and the missing estimates}

For any $C^1$ planar velocity on a buffer containing its trajectories,
let $\mathcal F_t$ be the flow from $t_0$ and put
\[
 K(t)=\int_{t_0}^t\|D_{x,y}u(s,\cdot)\|_{\mathrm{op},\infty}\,ds.
\]
The variational equation and its integral form give
$\|D\mathcal F_t\|\le e^K$ and
$\|D\mathcal F_t-I\|\le e^K-1$.
For a flat initial graph its tangent therefore has horizontal component
at least $2-e^K$ and vertical component at most $e^K-1$.  If
$K<\log2$, it remains a graph on its advected height interval and
\begin{equation}
 \|\epsilon_T\|_\infty\le\frac{e^K-1}{2-e^K}.
 \label{eq:cc-strain-bound}
\end{equation}
This sufficient bound applies to $p$ on a nonvanishing buffer and to
$u_\rho$ at fixed parameters.  It is not a canonical strain estimate.
The exact graph slope $r=\epsilon_T$, along $\dot T=u_1$, satisfies
\[
 \dot r=\partial_xu_2+(\partial_yu_2-\partial_xu_1)r
                         -(\partial_yu_1)r^2,
\]
so actual tangent-flow estimates can retain information discarded by
the norm in $K$.

For the scalar model, on a common height band covered by the graphs,
a concrete additional input is the moving-trace quantity
\[
 \mathcal B_\rho=\int_{t_0}^{t_1}\int_A^B
 (1+\epsilon_T^2)k_\rho(H)|H_{zz}|^2\,dT\,dt.
\]
Since $|j_\rho|^2=\mu_\rho k_\rho(H)|H_{zz}|^2$ and
$\mu_\rho/(\lambda+\mu_\rho)^2\le1/(4\lambda)$,
\begin{equation}
 \int\!\int|\epsilon_t|^2\le\frac{\mathcal B_\rho}{4\lambda},
 \qquad
 \frac1{B-A}\int_A^B|\epsilon(T,t_1)-\epsilon(T,t_0)|^2\,dT
 \le\frac{(t_1-t_0)\mathcal B_\rho}{4\lambda(B-A)}.
 \label{eq:cc-trace-bound}
\end{equation}
These follow from \eqref{eq:cc-compliant-flow} and Cauchy--Schwarz.
An independent estimate for $\mathcal B_\rho$ must retain its slope,
small-value weight and adaptive trace.  A fixed raw-field regulator can
freeze motion merely because the xi carrier decays.  If a prescribed
nonvanishing carrier $B$ is removed, $f=H/B$ requires instead the exact
current $-k_\rho(f)f_t\overline{f_z}$, with
\[
 -f_t=f_{zz}+2(B_z/B)f_z+(B_{zz}+B_t)f/B.
\]
Reflection-based bottom-flux cancellation also requires real-type $B$.
No uniform $\rho\downarrow0$ result for these adaptive boundaries is
asserted, nor may a variable regulator be substituted without its
derivative terms.  A stochastic approximation would likewise require
statistics along the actual moving trajectories: zero mean velocity
alone does not bound accumulated displacement or positive variation.

\subsection{The physical-jet connection and the topological distinction}

The original topological observable uses $G=(H,H_z)$, not $H\,dz$.
On $G\ne0$ put
\begin{equation}
 E=|H|^2+|H_z|^2,\qquad
 \alpha_G=\frac{\Im(G^\dagger dG)}{2\pi E},\qquad
 \mathcal C=\alpha_G\wedge d\alpha_G.
 \label{eq:cc-jet-connection}
\end{equation}
Here the full differential for $\mathcal C$ is in complex time--space.
Its integral over a complete regular three-dimensional boundary is the
normalized degree of $G/\|G\|$; isolated holomorphic joint zeros inside
contribute their positive local degrees, with the complex orientation
and the complete-boundary convention of
Section~\ref{sec:complete-contours}.  Regular-boundary homotopy
invariance permits interior events and does not select a contour
velocity.  A spatial line integral cannot replace this three-form.

On the real-time slice define
\[
 L_G=H_z^2-HH_{zz},\quad N_G=H_zH_{zz}-HH_{zzz},\qquad
 \mu_G=\frac{|L_G|^2}{\pi E^2},\quad
 j_G=\frac{N_G\overline{L_G}}{\pi E^2}.
\]
Then the exact curvature is
\begin{equation}
 d\alpha_G=\mu_Gdx\wedge dy+
             (\Im j_G)dt\wedge dx-(\Re j_G)dt\wedge dy.
 \label{eq:cc-jet-curvature}
\end{equation}
To verify it, in the chart $b=H_z\ne0$, $q=H/H_z$, write
\[
 \alpha_G=\frac1{2\pi}\left(d\arg b+
             \frac{\Im(\overline q\,dq)}{1+|q|^2}\right),\quad
 d\alpha_G=\frac{d\Re q\wedge d\Im q}{\pi(1+|q|^2)^2}.
\]
Substitute $q_z=L_G/b^2$ and $q_t=-N_G/b^2$.
The resulting polynomial ratios extend through $H_z=0$ wherever
$E>0$.  They are projective curvature and current, not scalar zero
density and current.

As before, for $\lambda\ge0$ where $\lambda+\mu_G>0$, the full
spacetime action with
\begin{equation}
 \mathcal A_G=\frac\lambda2(x\,dy-y\,dx)+\alpha_G
 \quad\hbox{gives}\quad
 u_G=\frac{j_G}{\lambda+\mu_G}.
 \label{eq:cc-jet-action}
\end{equation}
The proof is the same contraction and Cartan calculation used for
\eqref{eq:cc-scalar-action}.  For $\lambda=0$ and $L_G\ne0$ it gives
$u_G=N_G/L_G$, preserving $q$.  Positive $\lambda$ removes an
$L_G=0$ singularity where $E>0$, but not the joint-zero singularity
$E=0$.  Complete closed material loops conserve $\int\mathcal A_G$;
open arcs retain endpoint terms.  This generally noninteger circulation
is different from the complete Chern--Hopf degree.  For example, for
$H=z$ one has $\int_{|z|=R}\alpha_G=R^2/(1+R^2)$, whereas the
argument-principle count is one.  The scalar regularized circulation
may even remain conserved while actual zeros cross its boundary.
Neither conservation law removes the complete-boundary or excision
requirements of the physical topological calculation.

\paragraph{Concrete next analytic task.}
The primary choice is the heat current $u=p$, since its material form,
positive density, primitive and Noether action agree without an added
parameter.  On a prescribed compact initial contour, fix its actual
image under $\mathcal P(t_0,\cdot)$ and continue that inverse image.
The next result must bound its actual tangent flow and continue the
selected inverse branch, including its spatial buffer and escape
boundary, from complete canonical-source data on a stated height/time
scale. Separation from every global critical value is stronger than
necessary for a specified regular branch; a branch-local comparison
should retain the relevant domain instead.  In particular, \eqref{eq:cc-material-form} gives the exact relation
\[
 Z_a(t)=\frac{H(t_0,Z_0)Z_{0,a}}{H(t,Z(t))},\qquad
 \frac{d}{dt}\log Z_a=\partial_z(H_z/H)(t,Z(t)),
\]
with a continuous logarithm along a regular trajectory. Its phase controls
turning, and its modulus controls stretching. The actual
constants in \eqref{eq:cc-fourier}--\eqref{eq:cc-polynomial-invariants},
not just existence of the hierarchy, must enter any new restriction.
One should not assume a global nonzero level $\Im\mathcal P=c$ at
bounded width: at $t=0$,
$\Im\mathcal P(T+iy)=\int_0^y\Re H(0,T+iv)\,dv$ tends to zero for
bounded $y$ by the xi carrier's bounded-strip decay.  Using the actual
initial image avoids that incompatible default.
The scalar and physical-jet models provide comparisons preserving
different quantities.  For the scalar comparison, the unresolved
inputs are coupled tangent strain, end flux and one-sided relative
current, with any regulator limit proved separately.  A useful bound
must improve a direct estimate from the same canonical information;
the identities and collision controls above do not yet supply it.

\section{Research priorities after the finite comparison}
\label{sec:outlook}
The strongest completed structural result is the complete-source coefficient
equation together with its commuting finite propagator. The forest polynomial
turns it into a quantitative scalar-filter interface. The corrected
Riemann--Siegel reference is the most concrete theta-specific input now
available on the microscopic clock. These results fit together, but their
composition still contains genuine acquisition requirements.

\begin{center}
\begin{tabular}{>{\raggedright\arraybackslash}p{0.28\linewidth}>{\raggedright\arraybackslash}p{0.61\linewidth}}
\toprule
Result & Assumption or remaining input \\
\midrule
Forest quotient & Polynomial algebra; a paid denominator margin for an enclosure.\\
Coefficient transport & Complete population and a common nonzero wall; integrated complete-source moments and errors.\\
Positive-time multiplicity transfer & Complete real adjacent-root motion; RH in the stated canonical global application.\\
Microscopic corrected reference & Imported real-time source theorem, fixed truncation chart; whole-wall margin still unacquired at growing height.\\
Finite theta certificate & Complete-source interval errors on the specified box; no extension to unbounded height.\\
\bottomrule
\end{tabular}
\end{center}

\subsection{A bounded next acquisition with a decisive question}
A useful next test would fix one genuinely high-height chart before examining
its zeros. For example, put \(\mathcal L=\log(x/(4\pi))\) and prescribe
\(t=\mathcal L^{-3}\), with spatial scales proportional to
\(\mathcal L^{-1}\). With natural unfolding
\(\nu\sim\mathcal L/(4\pi)\), this gives \(4t\nu^2\to0\).
The choice is a proposed experimental design, not a theorem about its
resulting roots. A fixed Riemann--Siegel cutoff, or explicitly paid cutoff
changes, is required.

The test should acquire the corrected normalized reference on every side of
the declared wall, its absolute source error, the complete population count,
coefficient errors, and the coupled forest denominator/residual margin.
It should report the resulting filter bound alongside an equally informed
direct fixed-time calculation. The decisive question is which term fails:
the wall margin, the source-to-coefficient conversion, degree growth, the
forest denominator, or the reference filter itself. One successful chart is
a validation of the acquisition method, not a limiting law. A broad sweep is
unjustified unless it tests a specified uniform estimate or cost model.

\subsection{The arithmetic estimate cannot be replaced by conditioning}
The long-term target is \eqref{eq:global-spectral-interface}: a vanishing
normalized filter and vanishing omitted mass on a family of complete
populations. Better source errors make a reference trustworthy; they do not
make its filter small. Even excellent whole-wall conditioning can enclose an
interior multiple zero. An independent arithmetic estimate must ultimately
control the relevant clustering of the complete reference, or the unchanged
field itself, on the prescribed joint height/time scale.

An averaged one-sided small-value estimate for the complete normalized field
and derivative is a possible direction, provided a proved consumer connects
it to the needed root counts and every interpolation, exceptional-set and
coverage loss is paid. This is a research suggestion, not an estimate supplied
here. Generic real-root heat dynamics, a mean zero count, or a statistical
model cannot stand in for the missing canonical input.

\subsection{Conditional roles for backward transport and curved paths}
A further backward calculation is justified when it can acquire the full
correlated residual in \eqref{eq:paid-duhamel}, including the explicit
reference heat defect, and improve the final certified error or cost relative
to direct acquisition from the same information. The finite propagator alone
does not establish that improvement. The scalar late-time implementation
with separately bounded algebraic errors has a different, unfavorable
amplification budget; its failure does not exclude a joint source estimate.

Complex time should be pursued only with a proposed complete-source error
bound that improves a specified path budget. Holomorphic continuation of the
reference is insufficient because the current imported field theorem is for
real time. A cap calculation must additionally retain every connector and
complement. The completed low-height range-quadrature pilot identifies a
spatial integration bottleneck, so further source precision or geometrical
rephrasing would not address its dominant error. Higher-order correlated
integration could be a bounded numerical study, with the same improvement
available to the direct baseline; it is a lower priority than the microscopic
arithmetic and conditioning questions.

The present stopping point is therefore a finite method with an explicit
source interface and an honest comparison, not an almost-complete proof of
simplicity. Density-one simplicity and every-zero simplicity remain distinct
targets; the latter cannot discard a sparse family of doubles. An RH-free
application must also establish endpoint reality independently.


\section*{Acknowledgements and verification status}
This working manuscript was prepared for Edward Baker with substantial
assistance from OpenAI Codex using \texttt{gpt-6-astra}, reasoning effort
\emph{ultra}. Separate agents using the same model checked derivations,
source formulas, interval implementations and stated scope. These checks
are not independent specialist refereeing, cross-model validation, or
formal proof verification. The Vandermonde and spectral-filter exposition
was subsequently clarified with OpenAI Codex using \texttt{gpt-6.1-sol},
reasoning effort \texttt{xhigh}. The draft is intended for the author's
mathematical review before any circulation or attribution as an approved
paper.

The conservation section develops brainstorm Note 16 and its same-model
Review 07; it remains in this working draft pending a substantive analytic
result suitable for the selective book. Its manuscript integration and
proof obligations are recorded in brainstorm Review 08.

Research Notes 07--15, their reviews and small execution records are
maintained in the parent investigation directory. The earlier analytical
controls include 1,333 counted exact-rational assertions, with separate
floating diagnostics. The complete-source certificate uses Arb interval
arithmetic through python-flint, as described in its section; numerical
checks support the stated finite claims only. Sources, receipts and
provenance are kept separately so the mathematical assertions do not
depend on a file hash.

\begin{thebibliography}{99}
\bibitem{polymath}
D.~H.~J. Polymath,
\emph{Effective approximation of heat flow evolution of the Riemann
\(\xi\) function, and a new upper bound for the de Bruijn--Newman constant},
arXiv:1904.12438v2 (2019).
\url{https://arxiv.org/abs/1904.12438}.

\bibitem{rodgers-tao}
B.~Rodgers and T.~Tao,
\emph{The de Bruijn--Newman constant is non-negative},
Forum of Mathematics, Pi 8 (2020), e6.
\url{https://arxiv.org/abs/1801.05914}.

\bibitem{arb}
F.~Johansson,
\emph{Arb: efficient arbitrary-precision midpoint-radius interval arithmetic},
IEEE Transactions on Computers 66 (2017), 1281--1292.
\url{https://arxiv.org/abs/1611.02831}.

\bibitem{noether}
E.~Noether,
\emph{Invariant Variation Problems},
original in Nachrichten von der Gesellschaft der Wissenschaften zu
G\"ottingen (1918), 235--257; English translation by M.~A. Tavel (1971),
arXiv:physics/0503066v3.
\url{https://arxiv.org/abs/physics/0503066v3}.

\bibitem{cole}
J.~D. Cole,
\emph{On a quasi-linear parabolic equation occurring in aerodynamics},
Quarterly of Applied Mathematics 9 (1951), 225--236.
\url{https://doi.org/10.1090/qam/42889}.

\bibitem{hopf}
E.~Hopf,
\emph{The partial differential equation \(u_t+uu_x=\mu u_{xx}\)},
Communications on Pure and Applied Mathematics 3 (1950), 201--230.
\url{https://doi.org/10.1002/cpa.3160030302}.

\bibitem{ibragimov2007}
N.~H. Ibragimov,
\emph{A new conservation theorem},
Journal of Mathematical Analysis and Applications 333 (2007), 311--328.
\url{https://doi.org/10.1016/j.jmaa.2006.10.078}.

\bibitem{ibragimov2006}
N.~H. Ibragimov,
\emph{The answer to the question put to me by L.V. Ovsyannikov 33 years ago},
Archives of ALGA 3 (2006), 53--80.
The complete volume is available from Blekinge Institute of Technology:
\url{https://bth.diva-portal.org/smash/get/diva2%3A835112/FULLTEXT01.pdf}.

\bibitem{anco}
S.~C. Anco,
\emph{On the incompleteness of Ibragimov's conservation law theorem and
its equivalence to a standard formula using symmetries and adjoint-symmetries},
Symmetry 9 (2017), article 33; arXiv:1611.02330v3.
\url{https://arxiv.org/abs/1611.02330v3}.

\bibitem{papanicolaou}
V.~G. Papanicolaou, E.~Kallitsi and G.~Smyrlis,
\emph{Entire solutions for the heat equation},
Electronic Journal of Differential Equations (2021), no.~44, 1--25.
\url{https://ejde.math.txstate.edu/Volumes/2021/44/papanicolaou.pdf}.

\end{thebibliography}
\end{document}
